\chapter{Methodology}\label{chap:methods}

%% ============================================================
%%  Problem 1 — ICD Code Prediction Methodology (from Old Report)
%%  NOTE: This methodology was proposed but NOT fully executed
%%        due to high computational-infrastructure costs.
%% ============================================================

\section{Problem 1: ICD Code Prediction Methodology (Proposed)}\label{sec:p1:methodology}

\emph{Note:} The methodology described in this section was proposed as part of the ICD-10-CM code prediction research (Problem~1). Due to high computational-infrastructure costs, the full pipeline was not executed; only preliminary baseline experiments and data exploration were completed (see Chapter~\ref{chap:results} for partial results). The framework is documented here for completeness and to guide future work.

\subsection{Introduction and Scope}\label{sec:p1:meth_intro}
This section sets forth a multi-phase methodology aimed at refining automated ICD-10-CM code prediction, with particular emphasis on rare codes. The proposed framework features an ontology-based disease grouping strategy, a pipeline for synthetic discharge summary generation, a hybrid model training approach that integrates both real and synthetic data, and an evaluation scheme designed to account for rare code distribution imbalances. The methodology follows a modular and extensible workflow so that it may be adapted to alternative scenarios involving imbalanced classification. Validation was planned on a real-world clinical dataset to assess the generalizability and utility of this framework.

\subsection{Ontology-Based Disease Grouping}\label{sec:p1:ontology_grouping}
Rare ICD-10-CM codes frequently exhibit extremely low prevalence, which hinders the ability of predictive models to learn discriminative features. To mitigate this challenge, an ontology-based grouping approach was envisioned. This process entails parsing established medical ontologies---such as SNOMED CT, UMLS, and the Orphanet Rare Disease Ontology---to merge hierarchical and associative relationships.

Each rare disease would be anchored in a knowledge graph, and closely related concepts would be identified via direct graph traversal. These items would then be clustered into a unified disease family, thereby providing a contextual lens through which the rare disease can be understood. By associating each rare disease with more common conditions, the model would receive exposure to realistic clinical context, improving its capacity to recognize the linguistic cues relevant to rare diagnoses. Additionally, ontology mapping ensures that no orphaned codes are disregarded, given that each disease is matched at least to a more general ICD-10-CM code if an exact match proves elusive.

This ontology-centric approach is expected to enable more robust learning for rare codes, because they would be presented alongside clinically related diagnoses. By placing rare codes in medically coherent families, the methodology reduces the isolation that often occurs when codes appear only a handful of times in conventional datasets.

\subsection{Synthetic Discharge Summary Generation}\label{sec:p1:synthetic_generation}
Because rare codes remain infrequent in standard patient records, a synthetic data generation process was designed to amplify the presence of these underrepresented labels. A retrieval-augmented generation pipeline was planned, wherein medical facts, professional guidelines, and relevant textual fragments would be retrieved to guide a large language model in producing draft discharge summaries.

\begin{itemize}
    \item \textbf{Structured Prompting:} Placeholders would be embedded for sections such as Chief Complaint, History of Present Illness, Hospital Course, Medications, and Discharge Diagnoses.
    \item \textbf{Disease Clusters:} Rare diseases would be linked with related diagnoses identified through ontology-based grouping, ensuring that realistic co-occurrences appear in the synthetic text.
    \item \textbf{Multi-Agent Feedback Loop:}
    \begin{itemize}
        \item A draft discharge summary would be generated by the language model.
        \item A critic agent would then examine the draft for logical inconsistencies and missing details.
        \item Iterative refinements would be made until the note meets plausibility and consistency criteria.
    \end{itemize}
\end{itemize}

The final synthetic discharge summaries would need to provide sufficient evidence for the relevant rare ICD-10-CM codes, mirroring the style and structure of authentic clinical documentation. This balanced approach is expected to enrich the training dataset for rare codes, thereby reducing distributional skew and enhancing overall model recall for these infrequent conditions.

\subsection{Deep Learning Model Training and Testing}\label{sec:p1:dl_training}
An augmented training corpus composed of both real and synthetic discharge summaries would subsequently be employed in training two deep learning architectures: (i) a transformer-based model, and (ii) a CNN-based model. The transformer-based model is expected to capture long-range dependencies, whereas the CNN-based model should excel at identifying localized textual patterns pertinent to specific codes.

A multi-label classification scheme would be used, assigning probabilities to all potential ICD-10-CM codes. To address class imbalance, strategies such as weighted sampling, focal loss, and class-weighted loss were planned. Synthetic summaries would be oversampled or weighted more heavily, ensuring that rare codes appear frequently enough in training batches. The training procedure would commence with binary cross-entropy and moderate class weighting, followed by the incremental introduction of focal loss during later epochs, thus allowing the model to fine-tune its attention on the most difficult labels.

A development set containing both real and synthetic notes would guide hyperparameter tuning and early stopping. To confirm the practical value of this method, the final evaluation would be conducted on a held-out portion of real MIMIC-IV discharge summaries, assessing whether improvements in rare code classification carry over to genuine patient data.

\subsection{Evaluation Metrics and Analysis}\label{sec:p1:eval_metrics}
The planned evaluation seeks to capture various dimensions of performance, with special focus on rare code detection. Several metrics would be reported, each addressing a different aspect of classifier efficacy. The following formulas present the principal metrics in detail (where $TP$ = true positives, $FP$ = false positives, $FN$ = false negatives, and $C$ = total number of classes):

\begin{enumerate}
    \item \textbf{Micro-F1}
    \begin{equation}\label{eq:p1:micro_f1}
        \text{Micro-F1} = 2 \times 
        \frac{\sum_{i=1}^{C} TP_{i}}
        {\sum_{i=1}^{C} (2 \cdot TP_{i} + FP_{i} + FN_{i})}
    \end{equation}
    Micro-F1 aggregates the contributions of all codes to compute the overall precision and recall across the entire label space.

    \item \textbf{Macro-F1}
    \begin{equation}\label{eq:p1:macro_f1}
        \text{Macro-F1} = \frac{1}{C} \sum_{i=1}^{C} \text{F1}_{i}
    \end{equation}
    Each class-level F1 score $\text{F1}_{i}$ is calculated, and their average is taken to highlight how the model performs on each individual label, giving rare codes equal weight to common codes.

    \item \textbf{Recall@\emph{k}}
    \begin{equation}\label{eq:p1:recall_at_k}
        \text{Recall@}k = 
        \frac{|\{\text{correct code in top-}k\}|}
        {N_{\text{samples}}}
    \end{equation}
    This metric indicates the fraction of cases in which the true codes appear among the top $k$ predictions, a setting relevant in coding-assistance systems that present a limited list of suggestions.

    \item \textbf{Exact Match Ratio}
    \begin{equation}\label{eq:p1:exact_match}
        \text{Exact Match} = 
        \frac{|\{\text{notes with all codes correct}\}|}
        {N_{\text{notes}}}
    \end{equation}
    This stricter criterion gauges the extent to which all predicted codes align perfectly with the ground truth set for each note.

    \item \textbf{Never-Predicted Rate}
    \begin{equation}\label{eq:p1:never_predicted}
        \text{Never-Predicted Rate} = 
        \frac{|\{\text{codes never predicted}\}|}{C}
    \end{equation}
    This measure highlights whether any codes---particularly rare ones---are completely ignored.
\end{enumerate}

Emphasis would be placed on determining whether the ontology-based grouping and synthetic data augmentations produce measurable gains in macro-F1 and reduce the never-predicted rate. A targeted error analysis would be conducted for codes that remain unrecalled, providing directions for further refinements.

\subsection{Processing MIMIC-IV as the Test Dataset}\label{sec:p1:processing_mimic_iv}

In this work, the MIMIC-IV database will be employed as the principal test dataset for automated ICD coding. This dataset contains de-identified clinical information collected from the Beth Israel Deaconess Medical Center (BIDMC), with particular emphasis on structured hospital records and unstructured discharge summaries. The following steps describe how the data will be prepared, managed, and validated in a manner consistent with privacy regulations and best practices in clinical data science.

\subsubsection{Data Extraction and Integration}\label{sec:p1:data_extraction}
The relevant tables from MIMIC-IV (e.g., \texttt{admissions}, \texttt{patients}, \texttt{diagnoses\_icd}, and \texttt{discharge}) will be identified and retrieved. These tables will then be joined on keys such as \texttt{subject\_id} and \texttt{hadm\_id}, thus creating a unified view of patient admissions, annotated diagnoses, and the corresponding discharge narratives. During this process, all admissions lacking a corresponding discharge summary or valid ICD-10 code assignments will be flagged for potential exclusion, ensuring that the resulting dataset remains suitable for supervised learning.

\subsubsection{Text Standardization and Cleaning}\label{sec:p1:text_cleaning}
Each discharge summary will be examined to remove any remaining placeholders or noise introduced by the de-identification process. Lowercasing and tokenization will be conducted if deemed beneficial for subsequent analyses, and the presence of medical abbreviations or numeric identifiers will be noted. Segmentation into sections (such as ``History of Present Illness'' and ``Medications'') may be carried out to reduce text complexity, though the entire narrative will remain intact if full-length modeling is desired.

\subsubsection{Handling Duplicates and Missing Data}\label{sec:p1:duplicates_missing}
Instances in which multiple, nearly identical discharge summaries exist for a single admission will be managed by detecting duplicates through textual similarity checks. These checks will reduce the risk of overrepresenting particular patients or hospital stays. Any admissions without sufficiently long or clearly documented notes will be handled through either imputation strategies or careful removal, thereby minimizing noise in model training and evaluation.

\subsubsection{Code Distribution and Rare Labels}\label{sec:p1:code_distribution}
A preliminary examination of the ICD-10 codes assigned to each admission will be performed to understand their frequency distribution. Since ICD codes often follow a long-tail pattern, additional attention will be devoted to codes that appear infrequently. If these rare labels contribute significantly to the overall error rate, strategies such as targeted oversampling or retrieval-augmented data enrichment may be considered, provided such techniques respect the constraints of de-identification and data plausibility.

\subsubsection{Database Management and Version Control}\label{sec:p1:db_management}
All data transformations and table joins will be executed within a controlled environment (e.g., DuckDB) to allow for efficient querying and seamless integration with Python-based analytic workflows. The resulting datasets, along with any derived tables, will be tracked using version control to ensure that incremental changes are systematically logged and reproducible. This practice will enable transparent experimentation and simplify the process of rolling back to earlier dataset states if necessary.

\subsubsection{Security and Compliance}\label{sec:p1:security_compliance}
All patient information in MIMIC-IV is protected via de-identification, and no direct identifiers (e.g., names, addresses) are present. It will be confirmed that the data use agreement from PhysioNet is upheld, particularly regarding re-identification risks and data sharing restrictions. Any synthetic expansions or textual outputs generated by the pipeline will be carefully reviewed to ensure that confidential details are not reintroduced.

\medskip
\noindent\emph{As noted above, the full execution of this proposed methodology was not completed due to high computational-infrastructure costs. Preliminary baseline experiments and data exploration results are presented in Chapter~\ref{chap:results}.}


%% ============================================================
%%  Problem 2 — SA-AKI Mortality Prediction Methodology
%%  PRESERVED VERBATIM from the Current Report
%% ============================================================

\section{Problem 2: SA-AKI Mortality Prediction Methodology}\label{sec:p2:methodology}

\subsection{Overview}
This section details the methodological pipeline that was designed to develop and validate predictive models for SA--AKI mortality. The workflow begins with the construction of a high-fidelity cohort from an ontology-integrated data warehouse and proceeds through three analytical stages: (i) exploratory data analysis for phenotyping, (ii) development of a supervised binary classifier for mortality (the primary task), and (iii) time-to-event survival analysis. A strict, one-time stratified 80/20 data split underpins the entire process to prevent information leakage. All data-dependent transformations were fitted on the training set only and subsequently applied to the test set for a single final evaluation. The development process was guided by the TRIPOD statement for transparent reporting.

\subsection{Data Sources, ETL, and Cohort (T0--T24)}
\textbf{ETL and Integration.} The analysis asset was built by ingesting and integrating multiple sources, primarily MIMIC--IV, with medical ontologies (RxNorm, ATC, SNOMED~CT) linked via the Unified Medical Language System (UMLS). This was managed in a PostgreSQL database, enabling reproducible, ingredient-level value sets for concepts like systemic antibiotics and vasopressors. A validation suite ensured unit harmonization (e.g., F$\rightarrow$C, lb$\rightarrow$kg), referential integrity, and logical consistency.

\textbf{Cohort Definition.} The study population was defined as adults ($\geq$18\,y) on their first ICU stay with a length-of-stay $\geq$24\,h. The SA--AKI cohort was identified using established clinical criteria:
\begin{itemize}
    \item \textbf{Sepsis (Sepsis--3):} Defined by suspected infection (a body-fluid culture near the administration of a systemic antibiotic) and an acute increase in organ dysfunction ($\Delta\mathrm{SOFA}\geq 2$). The sepsis onset time, $t_{\text{sepsis}}$, was the start of the first qualifying ICU day.
    \item \textbf{AKI (KDIGO 2012):} Defined by changes in serum creatinine (SCr) or urine output. A hierarchical method established baseline SCr. The AKI onset time, $t_{\text{aki}}$, was the earliest time any criterion was met.
    \item \textbf{SA--AKI:} Required AKI onset within the window $[t_{\text{sepsis}}-48\text{ h},\,t_{\text{sepsis}}+48\text{ h}]$.
\end{itemize}
Patients with prior End-Stage Renal Disease (ESRD) or who received Renal Replacement Therapy (RRT) within the first 24\,h were excluded. All predictive features were restricted to measurements within the first 24 hours of ICU admission (T0--T24) to ensure clinical actionability and prevent look-ahead bias.

\subsection{Data Split and Leakage Control}
A one-time, stratified 80/20 train-test split was performed on the cohort. Let $\mathcal{D}=\{(x_i,y_i)\}_{i=1}^n$ be the full dataset, where $x_i\in\mathbb{R}^p$ are the T0--T24 features and $y_i\in\{0,1\}$ is the in-hospital death indicator. The partition is
\[
\mathcal{D} = \mathcal{D}_{\text{train}}\;\dot\cup\;\mathcal{D}_{\text{test}},
\]
where partitions are disjoint and stratified by the outcome to maintain class balance. All data-dependent objects—imputation models, scaling parameters, and encoding vocabularies—were fitted \emph{only} on $\mathcal{D}_{\text{train}}$ and then applied as frozen transformers to $\mathcal{D}_{\text{test}}$. Any validation set used for model tuning or calibration was created as a further split \emph{from within} $\mathcal{D}_{\text{train}}$.

\textbf{Leakage Audit.} A scripted audit ensured that no variables encoding information after T24 were included as predictors. This included removing features like total ICU length of stay and aggregates (e.g., \texttt{\_n} observation counts) that could act as proxies for care intensity or illness duration.

\subsection{Preprocessing and Feature Engineering (Train-Only)}
\textbf{Missingness Analysis.} The nature of missing data was formally investigated on $\mathcal{D}_{\text{train}}$. Little's MCAR test was performed, and its rejection of the null hypothesis suggested that data were not Missing Completely at Random (MCAR). This justified a more careful handling of missingness beyond simple imputation.

\textbf{Missingness Indicators.} For every feature with missing values, a binary indicator variable (\texttt{\_missing}) was created. The association between each indicator and the mortality outcome was tested using a chi-square test. Only indicators that showed a statistically significant association ($p<0.05$) were retained and appended to the dataset as new features, thereby preserving potentially prognostic information contained in the pattern of missingness.

\textbf{Imputation.} After adding significant missingness indicators, remaining missing values in numerical columns were imputed using the median calculated from the training set.

\textbf{Feature Selection and Scaling.} Highly correlated features (Pearson correlation > 0.99) were identified, and one from each pair was removed to reduce multicollinearity. All numerical features were then standardized using \texttt{StandardScaler}, with the scaler being fit only on the training data.

\subsection{Exploratory Data Analysis and Phenotyping (Training Only)}
An extensive EDA was performed on $\mathcal{D}_{\text{train}}$ to characterize the cohort's clinical properties. This served as an unsupervised phenotyping step. Key analyses included:
\begin{itemize}
    \item \textbf{Distributional Analysis:} Visualizing the distributions of key continuous variables (e.g., age, APACHE III, SOFA score, creatinine, lactate) and comparing them between survivor and non-survivor groups using histograms and density plots.
    \item \textbf{Correlation Analysis:} Calculating the Pearson correlation of all features with the mortality outcome to identify the strongest univariate predictors.
    \item \textbf{Stratified Analysis:} Using violin plots, box plots, and heatmaps to compare vital signs, laboratory values, and SOFA sub-scores between outcomes. Statistical significance was assessed using the Mann-Whitney U test due to the non-normal distribution of most variables.
\end{itemize}

\subsection{Supervised Binary Mortality Prediction (Primary)}
\textbf{Learning Objective.} For a model $f_\theta:\mathbb{R}^p\!\to\![0,1]$, the goal was to minimize the average binary cross-entropy loss on $\mathcal{D}_{\text{train}}$ using 5-fold cross-validation.

\textbf{Model Benchmarking.} A wide set of models was benchmarked: Logistic Regression, LDA, QDA, Naive Bayes, regularized linear models, Decision Tree, Random Forest, XGBoost, and LightGBM. For tree-based ensembles, hyperparameters were tuned using Optuna. The impact of class imbalance was addressed primarily through \texttt{class\_weight='balanced'} and was compared against various oversampling strategies (e.g., SMOTE, ADASYN) in sensitivity analyses.

\textbf{Final Model Selection, Calibration, and Thresholding.} The LightGBM model was selected based on its superior performance. The final pipeline involved:
\begin{enumerate}
    \item Training the optimized LightGBM model on a sub-split of the training data.
    \item Calibrating the model's outputs on a held-out validation split (from the training set) using \texttt{CalibratedClassifierCV} with the isotonic method.
    \item Determining the optimal classification threshold on the same validation split by maximizing Youden's J statistic.
    \item Retraining the final, calibrated LightGBM model on the entire training dataset ($\mathcal{D}_{\text{train}}$).
    \item Evaluating the final model on the untouched test set ($\mathcal{D}_{\text{test}}$) using the predetermined threshold.
\end{enumerate}

\textbf{Interpretability.} The final LightGBM model was interpreted using SHAP. A beeswarm summary plot was generated to visualize global feature importance and the direction of effects.

\subsection{Survival Analysis (Time to Death)}
The same preprocessed T0--T24 feature matrix was used to model time to in-hospital death, with patients discharged alive being right-censored.

\begin{itemize}
    \item \textbf{Cox Proportional Hazards (PH):} A CoxPH model was fitted using the \texttt{lifelines} library. Feature selection was performed by retaining the top 50 predictors from an initial full model, and the penalizer was tuned. Results were interpreted via a forest plot of hazard ratios.
    \item \textbf{DeepSurv:} A custom neural network-based Cox model was implemented in PyTorch, using SELU activations and residual connections. The model was trained by minimizing the negative Cox partial log-likelihood with the Efron method for handling ties.
\end{itemize}
Model performance was assessed by Harrell's C-index, time-dependent AUC, calibration plots at clinically relevant horizons (90, 180, and 365 hours post-admission), and Kaplan-Meier curves stratified by model-predicted risk quartiles. The statistical significance of the risk stratification was confirmed with the log-rank test.

\subsection{Adherence to TRIPOD Guidelines}

\subsubsection{Data Leakage: Definition, Mechanisms, and Mitigation}
In the context of clinical prediction modeling, data leakage has been defined as the inadvertent incorporation of information that would not be available at the intended time of prediction into the model training or evaluation process, resulting in inflated performance estimates that fail to reflect real-world utility \citep{Kang_Yoon_2023}. For SA--AKI mortality prediction, three principal mechanisms through which leakage may be introduced have been identified in the literature and were considered during the design of this pipeline.

The first mechanism, temporal leakage (also termed look-ahead bias), arises when features are derived from measurements recorded after the intended prediction horizon. A model designed to estimate mortality risk at 24 hours post-admission that incorporates laboratory values from subsequent days, or whole-stay aggregates such as the mean creatinine across the entire ICU episode, has been provided with information that encodes the patient's future trajectory and outcome. The inclusion of total ICU length of stay as a predictor represents a particularly direct form of this bias, as length of stay serves as a strong proxy for whether and when death occurred \citep{Chen_2025}.

The second mechanism, preprocessing leakage, has been observed when data-dependent transformations (including imputation, scaling, and feature selection) are fitted on the entire dataset prior to the train--test split. Under such conditions, the statistical properties of the held-out records (means, variances, and correlations with the outcome) are permitted to influence the training pipeline, and optimistically biased performance estimates are produced as a consequence \citep{Kang_Yoon_2023}.

The third mechanism, selection leakage, occurs when feature screening is conducted across the full dataset rather than within the training partition alone. Features that happen to correlate with the outcome across all records, including those subsequently reserved for evaluation, are preferentially retained, and the resulting predictor set carries information from the evaluation data.

In the present work, each of these mechanisms has been addressed through explicit safeguards. Temporal leakage has been prevented by restricting all features to the T0--T24 window and by deploying a scripted audit that flags any variable requiring post-T24 information. Preprocessing leakage has been avoided by fitting all transformers (imputer, scaler, and missingness indicator selector) exclusively on $\mathcal{D}_{\text{train}}$. Selection leakage has been precluded by confining all feature screening to the training partition. The held-out test set was used exactly once for final evaluation, and no iterative refinement based on test-set performance was undertaken.

\subsubsection{The 24-Hour Window: Rationale, Strengths, and Heterogeneity Considerations}\label{sec:t24-heterogeneity}
The selection of a fixed 24-hour observation window anchored to ICU admission (T0--T24) was motivated by the structure of clinical workflows in intensive care. The first ICU day is the period during which the initial assessment is typically completed, baseline investigations become available, and early management decisions regarding fluid resuscitation, vasopressor initiation, and antibiotic escalation are made. A prediction generated at T24 is therefore aligned with the natural decision point at which escalation or de-escalation of care is considered.

A fixed clock anchored to ICU admission, however, introduces a source of heterogeneity that warrants explicit discussion. Patients arrive in the ICU at different stages of their illness trajectory: some are admitted directly from the emergency department near the onset of sepsis, whereas others are transferred from general wards after a period of progressive deterioration. Consequently, the T0--T24 window captures different phases of the underlying disease process across patients. For early presenters, the window may coincide with the initial inflammatory surge, while for late transfers, it may reflect a more established pattern of organ dysfunction.

This heterogeneity represents a genuine limitation of any admission-anchored design, yet it also reflects the conditions under which a bedside prediction tool would be required to operate in practice. A clinician querying the model at the 24-hour mark has access only to the information accumulated since admission, regardless of the patient's prior trajectory. Several features included in the model were observed to partially account for this variation: the APACHE~III score, computed at admission, integrates acute physiology with chronic health status and thereby captures elements of the pre-ICU trajectory; the SOFA score at 24 hours reflects cumulative organ dysfunction irrespective of when it began; the Charlson comorbidity index components provide context regarding baseline vulnerability; and the first/last/delta/slope suffixes applied to time-varying features encode within-window dynamics that have been expected to differ between early and late presenters.

Stratified analyses by admission source (emergency department versus ward transfer) or by the interval from symptom onset to ICU admission could be pursued in future work to quantify the degree to which model performance varies across these subgroups. An alternative design in which the observation window is anchored to sepsis onset rather than ICU admission has also been considered, though such an approach introduces its own difficulties because the sepsis onset time is itself an estimated quantity derived from clinical rules rather than a directly observed event.

\medskip
The present study was designed in accordance with the TRIPOD (Transparent Reporting of a multivariable prediction model for Individual Prognosis Or Diagnosis) guidelines. The key items addressed include:
\begin{itemize}
    \item \textbf{Source of Data (Item 4a):} The study design and source of data (MIMIC-IV, a large electronic health record database) are clearly described.
    \item \textbf{Participants (Item 5a-b):} The study setting (ICU), location, and detailed eligibility criteria for participants are specified in the cohort construction.
    \item \textbf{Outcome and Predictors (Items 6a, 7a):} The outcome (in-hospital mortality) and all predictors (derived from the T0--T24 window) are explicitly defined.
    \item \textbf{Sample Size (Item 8):} The flow of participants through the study is described, detailing exclusions and the final sample size available for analysis.
    \item \textbf{Missing Data (Item 9):} The handling of missing data is detailed, including statistical testing of the missingness mechanism, creation and selection of significant missingness indicators, and median imputation based on the training set.
    \item \textbf{Statistical Analysis Methods (Item 10a-d):} Predictor handling (correlation removal, scaling), model-building procedures (benchmarking of multiple algorithms, Optuna for tuning), methods for internal validation (stratified 80/20 train-test split), and all performance measures (AUROC, F1-score, C-index, calibration plots, DCA) are specified.
    \item \textbf{Model Specification and Performance (Items 15a, 16):} The final model (LightGBM) and its key hyperparameters are presented. Performance measures on the test set are reported, providing an unbiased estimate of the model's performance on new data.
\end{itemize}
